\chapter{Análisis de metodologías de desarrollo}
\label{cap:analisis_metodologías}

Las metodologías de desarrollo software han evoluciado con el tiempo para adaptarse a las diferentes necesidades 
y contextos de proyectos. En función de las características del proyecto y del equipo de desarrollo,
será más conveniente usar una u otra. 

En este capítulo se analizarán las metodologías de desarrollo de software más utilizadas en la actualidad. 
Se describirán las características de cada una de ellas, así como sus ventajas y desventajas. Además, se realizará una comparativa entre las metodologías más populares 
para determinar cuál es la más adecuada para el desarrollo de la aplicación.

\section{Metodologías de desarrollo de software}
\label{sec:metodologias}

Las metodologías de desarrollo de software son un conjunto de prácticas y técnicas que se utilizan para planificar, diseñar, implementar y probar un proyecto de software. A continuación,
se describen las metodologías más utilizadas en la actualidad:

\subsubsection{Modelo en cascada (Waterfall)}
\label{subsec:waterfall}

Una de las metodologías más tradicionales y conocidas es el \textbf{Modelo en Cascada} (Waterfall), el cual sigue un enfoque secuencial en el que cada fase del 
desarrollo debe completarse antes de avanzar a la siguiente. Esta metodología se compone principalmente de las estapas de recopilación de requisitos, diseño, implementación, pruebas, despliegue y mantenimiento. 
Su principal ventaja es la claridad de los requisitos desde el inicio, lo que facilita una buena planificación y documentación. Sin embargo, su principal inconveniente
es la falta de flexibilidad, ya que cualquier cambio en los requisitos implica volver a fases anteriores, lo que puede
generar retrasos y costos adicionales. Esta metodología sigue siedno utilizada en proyectos donde los requisitos 
están bien definidos y no se esperan cambios significativos. \cite{sommerville2015}


\subsubsection{Modelo en espiral}
\label{subsec:espiral}

El \textbf{Modelo en Espiral} introduce un enfoque iterativo que prioriza la gestión de riegos. En cada iteración se realiza actividades
de planificación, análisis de riegos, implementación y evaluación, lo que permite detectar y mitigar problemas antes de que se conviertan en fallos grandes. 
Esta metodología es útil en proyectos grandes y complejos, donde la incertidumbre es alta. No obstante su aplicación suele ser costosa y requiere un gestión constante del proceso,
lo que lo hace menos adecuado para proyectos pequeños. \cite{boehm1988}

\subsubsection{Desarrollo iterativo e Incremental}
\label{subsec:iterativo}
